Representing semantic knowledge as subject--predicate--object triples forces
higher-arity and context-scoped facts through reification: one statement
becomes many records, and answering a simple lookup becomes a multi-way
join. We describe SFDB, a prototype database whose organising principle is
instead the \emph{context poset}: each fact is stored atomically as a
section over a finite poset of contexts, and entity- and context-anchored
queries are
answered through precomputed indexes rather than reconstruction. The formal
vocabulary is that
of presheaves on a finite Alexandrov space; the sheaf condition holds
vacuously there, and we claim no category-theoretic novelty. A reified
triple-store baseline is co-implemented over the same canonical schema, and
every benchmark query is checked, online, for canonical result equivalence
across engines --- a harness that caught a real semantic bug in our own
engine before publication.

Across \benchscalecount{} scales (\benchscales{} facts), SFDB answers
LOOKUP queries \rangelookupkg{} faster, full-scan GLOBAL queries
\rangeglobalkg{} faster, and CONTEXT queries --- retrieval scoped to a
context and its descendants, the query shape the design is named for ---
\rangecontextkg{} faster than the baseline. A storage-layer control --- the
same baseline with its SQLite tables replaced by in-process dict indexes
--- shows that most of the LOOKUP, GLOBAL, and CONTEXT advantage
(\rangelookupkgmem{}, \rangeglobalkgmem{}, and \rangecontextkgmem{})
survives with the storage-layer difference removed
and is therefore attributable to the design itself. The
temporal-range advantages do not survive the control at all --- SFDB is at
parity or measurably slower than the control on bounded temporal ranges at
most scales --- and insert throughput is slower than the baseline at the
smallest scale measured and faster at every larger one; both are reported
as negative or mixed results rather than smoothed into a single headline.
The cost that funds the advantages that do hold is
memory: \rangememkg{} more resident memory per fact than the baseline.
Against three real, disk-backed production stores accessed over
SPARQL/HTTP --- Apache Jena TDB2, OpenLink Virtuoso, and Ontotext
GraphDB, independently
engineered by three different organisations --- the pattern holds up on LOOKUP,
GLOBAL, and CONTEXT --- SFDB remains faster than all three at every scale, despite
their latencies including an HTTP round trip SFDB's in-process calls do
not pay --- but reverses on TEMPORAL from $10\,000$ facts
onward against all three stores, where a mature production optimiser
overtakes SFDB's Python-level candidate filtering. We report this
crossover as measured rather than smoothing it into either a clean win
or a hedge; that it appears independently against three unrelated
optimisers is evidence it is a property of mature query engines
generally rather than one vendor's implementation. Two further
comparisons --- \wikidatanumfacts{} genuine Wikidata facts we did not
construct, and LUBM, the standard RDF benchmark workload, at four scales
up to \lubmlargescale{} facts --- extend the finding rather than merely
repeat it: LOOKUP, GLOBAL, and CONTEXT reproduce cleanly on both, while
TEMPORAL lands in a different place on each (beating both in-house
baselines on Wikidata, roughly at parity with them on LUBM), direct
evidence that only TEMPORAL among the four query classes is genuinely
sensitive to a dataset's actual shape. Building the LUBM comparison
surfaced the evaluation's most consequential self-correction: a
performance bug in SFDB's own CONTEXT query resolution, present for the
entire evaluation and exposed only by a topology shape (many large,
mutually unrelated top-level contexts) none of this paper's other
datasets have, which required re-running and re-verifying every
CONTEXT number this paper reports and retracting two earlier claims
built on the unfixed numbers. Four external reference points (rdflib, Jena TDB2, Virtuoso, and
GraphDB,
all five query classes, result counts verified against all four), two
real-world/standard-workload comparisons, and ten documented
rounds of self-correction --- including the CONTEXT bug above and a
cross-engine SPARQL semantics divergence found while building the
Virtuoso comparison ---
complete an evaluation built to be checked rather than believed. All
results regenerate from the released artifact with a single command.
